% 本文件由 LaTeX 排版 Agent 根据有效 translation.json 自动生成。
% author=Ephraim the Syrian; work_id=3706; title=论我们的主

\chapter{论我们的主}

% structure: p0001 source=h1 target=omitted reason=page-title-already-rendered-by-parent

因此，愿众口颂扬那从他们身上除去亵渎言论的祂。光荣归于祢，祢离开一个居所，去寄居在另一个居所！为使祂来临，并使我们成为派遣祂者的居所，\OriginalTerm{独生子}{only-begotten}离开【与】\OriginalTerm{神性}{Deity}【同在】，寄居在童贞玛利亚内；这样，借着普通的出生方式，虽是独生子，祂却成为许多人的兄弟。祂离开了\OriginalTerm{阴府}{Sheol}，寄居在\OriginalTerm{天国}{Kingdom}中；为要从压迫众人的阴府，寻找一条通往赏报众人的\OriginalTerm{天国}{Kingdom}的道路。因为我们的主将自己的复活作为抵押给了必死之人，为要将他们从毫无区别地接纳逝者的阴府，迁移到有区别地接纳被邀请者的\OriginalTerm{天国}{Kingdom}；好使我们从【那】使众人身体在其中平等的【计划】，来到【那】在其中区分众人行为的【计划】。祂就是那位下到阴府又上升的，为要将我们从那腐蚀寄居者的地方，带到那以福乐滋养其居民的地方；甚至那些居民，他们用这个世界的财物、果实和花朵——这些都要逝去——为自己在那里装饰了永不逝去的帐幕。那按本性受生的首生子，在祂本性之外的另一次出生中诞生，为使我们知道，在我们自然的出生之后，我们必须有另一次在我们本性之外的出生。因为祂既是属神的，直到祂来到肉身的出生之前，不能成为有肉身；同样，属肉身的，除非他们在另一次出生中出生，不能成为属神的。但那生育不可探究的圣子，在另一次可探究的出生中诞生；为使我们从前者学到祂的尊威没有界限，从后者学到祂的恩宠没有量度。因为祂的尊威伟大无量，祂的第一次受生在我们的任何思想中都无法想象。祂的恩宠丰盈无限，祂的第二次诞生被众口传扬。\par

2. 这是那一位按祂的本性由\OriginalTerm{神性}{Godhead}所生，又按不是祂本性的方式由\OriginalTerm{人性}{manhood}所生，并按照不是祂惯例的方式由圣洗所生；好使我们按我们的本性由人性所生，又按不是我们本性的方式由神性所生，并按照不是我们的惯例借着圣神所生。祂既由神性所生，就来到了第二次出生；为将我们带到那被谈论的出生，就是祂由圣父所生——不是要探究，而是要相信——以及祂由女人所生，不是要轻视，而是要高举。现在祂在十字架上的死亡见证了祂由女人所生。因为那死了的，也曾出生。加俾额尔的\OriginalTerm{圣母领报}{Annunciation}宣告了祂由圣父所生，即\BibleQuote{至高者的能力要庇荫祢} \BibleRef{路 1:35}。如果那真是至高者的能力，显然就不是凡人的种子。所以祂在母胎中的受孕与祂在十字架上的死亡相连；祂第一次的出生与天使的宣告相连；为叫那否认祂出生的人被祂的钉死所驳斥，而那以为祂始于玛利亚的人，得以受劝诫：祂的神性在万有之先；所以那断定祂的起始为肉身的，【可因叙述了祂由圣父发出的记载而被证明有误】。圣父生了祂，并借着祂创造了受造物。肉身生了祂，并借着祂歼灭了情欲。圣洗产生了祂，为借着祂洗净污渍。阴府产生了祂，为借着祂掏空其宝库。祂从圣父身边来到我们这里，经过受生者的道路；又经过死亡者的道路，前往圣父那里；好使祂借着出生而来临，祂的来临得以被看见；祂借着复活而返回，祂的离去得以被确认。\par

3. 但我们的主被死亡践踏；祂反过来开辟了一条越过死亡的道路。祂就是那位甘愿服从并忍受死亡的主，为要违背死亡的意愿将其推翻。因为我们的主背着十字架前行，顺从死亡的意愿；但祂在十字架上呼喊\BibleRef{玛 27:50}，便违背死亡的意愿，从阴府内带出了死者。因为死亡用来杀害祂的那同一事物（即身体），祂以此为武器取得了对死亡的胜利。但\OriginalTerm{神性}{Godhead}隐藏在人性中与死亡作战；死亡杀了，也被杀了。死亡杀了自然生命，超自然生命却杀了死亡。由于死亡不能没有身体而吞噬祂，阴府也不能没有血肉而吞没祂，祂来到童贞女那里，为要从那里取得那能将祂带往阴府的形体；正如人们从驴旁为祂牵来驴驹，祂骑着它进入耶路撒冷，并宣告那隐藏的覆灭和其子女的毁灭。于是，带着来自童贞女的身体，祂进入阴府，洗劫了它的库房，清空了它的宝藏。祂来到众生之母厄娃那里。这葡萄树，死亡借着她自己的手拆毁了它的篱笆，使她尝到了它的果实。这样，众生之母厄娃成了所有活着的人的死亡之源。但玛利亚萌芽了，一棵从古葡萄树厄娃发出的新枝，新生命居住在她内，为的是当死亡习惯性地来吞噬可朽的果实时，那杀死死亡的生命得以隐藏在其中对抗他；为的是当死亡毫不畏惧地吞下果实后，他可能连同它们一起吐出许多果实。因为\OriginalTerm{生命之药}{Medicine of life}从天降下，与身体——那可朽的果实——融合。当死亡照例来吞噬时，生命反过来吞噬了死亡。这食物饥饿地想要吃掉它的吞食者。因此，通过死亡贪婪地吞下的一个果实，他吐出了他贪婪地吞下的许多生命。那催逼他攻击一个的饥饿，倒空了他催逼他攻击许多的贪婪。这样，死亡忙于吞下一个，却急于释放许多。因为当一位在十字架上死去时，许多埋葬在阴府里的人因祂的呼喊而出来。\BibleRef{玛 27:50}这就是那果实，它劈开了吞噬它的死亡，从其中带出了那被派遣去追寻的生命。因为阴府隐藏了她所吞噬的一切。但通过那位未被吞噬的一位，一切被她吞噬的都从她内被恢复。一个人的胃紊乱了，他会吐出既甜又不甜的东西。同样，死亡的胃紊乱了，当它吐出那使其生病的生命之药时，它也一同吐出了那些曾被它愉快地吞噬的生命。\par

4. 这就是木匠之子，祂巧妙地使祂的十字架成为一座跨越吞噬一切的阴府的桥梁，将人类带入生命的寓所。因为人类正是通过树而落入阴府，所以也是通过树他们进入了生命的寓所。因此，通过那棵树，在其中尝到了苦味，也通过它尝到了甜味，为叫我们知道，在受造物中，没有什么能抵抗祂。光荣归于祢，祢将祢的十字架架设成一座跨越死亡的桥梁，好使灵魂能借着它从死者的寓所渡到生命的寓所！\par

5. 外邦人赞美祢，因为祢的\OriginalTerm{圣言}{Word}在他们面前成为一面镜子，使他们能在其中看见死亡秘密地吞噬他们的生命。但是偶像被工匠装饰着，而它们的装饰却毁损了它们的装饰者。然而祢将他们吸引到祢的十字架；当身体的美丽在十字架上被毁损时，心灵的美丽却在上面闪耀。于是，那些跟随不是神的神祇的外邦人，那一位真神去追他们，并藉祂的话语，如同藉缰绳，将他们从众神转向独一之神。这就是那位\OriginalTerm{大能者}{Mighty One}，祂的宣讲成为外邦人嘴中的缰绳，引导他们离开偶像走向那派遣祂者。但死寂的偶像，闭着嘴，靠崇拜者的生命为生。因此，祢将祢的血混入他们的肉中，这血削弱并推倒了死亡，好使他们的吞噬者的嘴从他们的生命中被赶走。同样，因为\Person{以色列}杀害了祢，被祢的血所玷污，那嫁接在他身上的偶像崇拜因祢的血被从他身上赶走。因为他通过祢的血断奶脱离了异教；在此之前，他从未断奶。\par

6. 但是以色列人将我们的主钉在十字架上，借口说祂确实引诱我们离开独一天主。但他们自己却常因许多偶像而离开独一天主。因此，当他们以为钉死那引诱他们离开独一天主的那位时，却发现被祂从所有偶像中引领到独一天主那里；为的是，因为他们没有自愿认识祂是天主，他们就不得不被迫认识祂是天主；时机一到，那藉着祂所赐予他们的恩惠，就要控告他们亲手所行的邪恶。这样，尽管压迫者的舌头否认，但他们所获得的帮助却定了他们的罪。因为恩宠超乎他们所能负担地压在他们身上，好使他们满载祢的祝福时，羞愧而不敢否认祢的位格。祢也怜悯了那些其生命已被献给死偶像作食物的人。他们在旷野所造的那只牛犊\BibleRef{出 32:4}，如同在旷野吃草一样，吞食了他们的生命。因为他们从埃及偷来并藏在心中的偶像崇拜，一旦显露出来，就公开杀死了那些秘密藏匿它的人。这好比藏在木头里的火，当它从内部生出时，就焚烧木头。梅瑟将牛犊磨成粉末，令他们喝下这试炼水\BibleRef{出 32:20}；好叫凡为敬拜牛犊而活的人，因饮下牛犊而死亡。肋未的子孙就奔向他们，那些奔向梅瑟并佩上刀剑的人。因为肋未的子孙不知道应当击杀谁，因为敬拜者与不敬拜者混杂在一起。但那容易分辨的主，却区分了玷污与未玷污的人；好让无辜者感谢他们的无辜没有在公义者面前被忽视，而有罪者被证实他们的过犯没有逃脱审判者的眼目。但肋未的子孙是公开的报应者。因此梅瑟在犯罪者身上做了记号，使报应者容易报应。因为牛犊的饮料进入了那些内心存有对牛犊之爱的人，并在他们身上显出明显的标记，使拔出的刀向他们冲来。于是那在牛犊崇拜中犯了奸淫的会众，他令他们喝下试炼水，好让淫妇的标记出现在其中。由此衍生出关于妇女的法律，即她们要喝试炼水\BibleRef{户 5:17}；好使淫妇身上出现的标记，提醒会众回念自己在牛犊崇拜中所犯的奸淫，并以敬畏之心防备另一种奸淫；同时以忏悔之心追念先前的奸淫；这样，当他们审判自己的妇女，若她们对他们不忠，他们就能定罪自己，因为他们对自己天主不忠。\par

7. 愿光荣归于祢，祢藉着祢的十字架除去了外邦人的习俗，这习俗曾使受割礼和未受割礼的人都失足跌倒！愿赞颂归于祢，生命的良药，祢使所有受洗者归向那万有的生命、万有的主！那被寻回的迷失者祝福祢；因为藉着寻回迷失者，祢使那些未被迷失的天使也得了喜乐。未受割礼者赞美祢，因为在祢的平安中，那中间存在的敌意被吞没了，因为祢接受了在祢肉身上的割礼外在标记，通过这标记，本来属于祢的未受割礼者常被算作不属于祢。因为祢以心中的割礼作为祢的标记；藉此，受割礼者被显明他们不属于祢。因为祢来到自己的地方\BibleRef{若 1:2}，自己的人却不接受祢；藉此他们被显明他们不属于祢。但那些祢没有来到的人，因祢的仁慈在祢后面呼喊，求祢用从孩子桌子上掉下来的碎屑满足他们。\par

8. 天主从神性中被派遣而来，以定罪那些雕像不是神。当祂从它们那里取走那装饰它们的天主之名时，它们本身的瑕疵就显现了。它们的瑕疵是——它们有眼却不能看，有耳却不能听。祢的宣讲说服了它们的众多敬拜者，使他们将众多神祇换为独一神。因为祢从偶像那里取走了神性的名称，敬拜也随名称一起被撤回；即那与名称相联的敬拜；因为敬拜也随从于天主的名称。因此，既然所有外邦人都向这名献上敬拜，最后这当受敬拜的名称将完全归于它的主。所以在最后的敬拜中，也将完全归于它的主，以应验\Emph{万物都要服从祂}。然后，祂也要把万物制服以后，\Emph{使自己服从那使万物服从祂的父}。\BibleRef{格前 15:27} 好让这名，从一级上升到另一级，与它的根源相连。因为当所有受造物因爱而联系于那创造他们的圣子，而圣子因爱联系于那生祂的圣父时，所有受造物最后都要感谢圣子，因他们藉着祂领受了一切恩惠；并在祂内并偕同祂，也要感谢祂的圣父，祂的宝库向我们分施一切财富。\par

9. 光荣归于祢，祢穿上了必死亚当的身体，并使之成为所有必死之人的生命泉源。祢是那永生的，因为祢的杀害者如同祢生命的农夫，他们将它像麦子一样播在地深处，好使它复活并一同复活许多人。来吧，让我们使我们的爱成为团体的巨大香炉，并在其上献上我们的颂歌和祈祷，献给那使祂的十字架成为神性的香炉、并为我们众人献上香祭的那位。那高高在上者俯就下方的人，以将祂的财富分施给他们。因此，虽然贫穷者接近祂的人性，却从神性领受恩赐。所以祂使祂所穿上的身体成为祂财富的掌管者，好使祂，主啊，从祢的宝库中取出，分施给贫穷者，祂亲族的子孙。\par

10. 愿光荣归于祂，祂从我们领受了，为要赐给我们；好使我们藉着那属于我们的，更丰盛地领受那属于祂的！是的，藉着那位中保，人类得以从它的助手中领受生命，正如起初藉着中保，人类从它的毁灭者那里领受了死亡。祢为自己造了身体作仆人，为要藉着它，将人们所渴望的一切赐予渴望祢的人。此外，在祢身上，那些杀害[祢]并埋葬[祢]的人隐藏的意愿得以显明；这是因为祢穿上了身体。因着祢那身体的缘故，杀害祢的人杀了祢，也死于祢手；因着祢身体的缘故，埋葬祢的人埋葬了祢，也与祢一同复活。那不可触摸的大能降下，穿上了可触摸的肢体，使贫乏的人能亲近祂，在触摸祂的人性时，他们能辨别祂的天主性。因为那哑巴（主所治愈的）用身体的指头感知到祂靠近了他的耳朵并触摸了他的舌头；\BibleRef{谷 7:32} 不，他以可触摸的指头触摸了不可触摸的天主性；那时祂正松解他舌头的捆绑，打开他耳朵堵塞的门户。因为身体的建筑师和肉体的工匠来到他那里，用温和的声音毫不痛苦地穿透了他变厚的耳朵。他那闭合的口，本不能生出言语，却生出了对祂的赞美，是祂使这不毛之地产出言语。那么，那赐给亚当不经教导就能立刻说话的祂，亲自赐给哑巴容易说话的能力——那些本难以学会的舌头。\par

11. 看，另一个问题又清楚了：——我们查考主以何种语言赐予那些从各方言语中来到祂那里的哑巴说话的能力？虽然这很容易知道，但我们的灵魂推动我们去追求那比这更大的知识。那知识便是：要知道第一个人类是藉着圣子造成的。因为从这一事实——即藉着祂，言语赐给了亚当的子孙哑巴——我们可以学到，言语是藉着祂赐给了他们的始祖亚当。在此，有缺陷的本性也被主补全了。那么，祂能补全本性的缺陷——显然，本性的补全是由祂制定的。但没有比这更大的缺陷了，即一个人生来不能说话。既然在于言语中我们超越所有受造物，它的缺陷便比所有其他缺陷更大。那么，祂，一切缺陷藉着祂得以补全——显然，一切圆满藉着祂得以建立。但因为藉着祂，所有肢体在母胎中秘密地领受圆满，所以藉着祂，他们的缺陷被公开补全；使我们学到，起初整个结构是藉着祂构成的。于是祂吐唾沫在指头上，放在那聋子的耳朵里；祂用唾沫和了泥，涂在那瞎子的眼睛上；\BibleRef{若 9:6} 使我们学到，正如那生来瞎眼的人的眼球有缺陷，这人的耳朵也有缺陷。因此，藉着来自那圆满者身体的酵，我们形成中的缺陷得以补全。因为主不应该从自己身上割下任何东西去补全其他身体的不足；而是用那可以从祂身上取走的东西，祂补全了那些缺乏之人的不足；正如在可以吃的东西中，有死之人吃祂。于是祂补全了缺陷，并赐予有死之人生命，使我们知道，从那充满圆满的身体中，缺乏之人的缺乏得以补全；从那充满生命的身体中，\BibleRef{哥 2:9} 生命赐予了有死之人。\par

12. 先知施行了所有其他奇事，但从未补全肢体的缺陷。但身体的缺陷被保留下来，要藉着主来补全；使灵魂能感知到，一切缺陷都必须藉着祂来补全。因此，有智慧的人应当觉察到，那补全受造物缺陷的祂，是创造者形成能力的掌管者。但当主在地上时，祂赐给聋哑人听力和说他们没有学过的语言的能力；使祂升天后，人能明白祂赐给祂的门徒说各种语言的能力。\par

13. 那些钉死祂的人以为，当我们的主死后，祂的征兆也随之消亡了。但祂的征兆却借着祂的门徒继续显明它们的存在，好让凶手知道征兆之主是活着的。事先，凶手们制造麻烦，喊着说祂的门徒偷走了祂的尸体。但后来，借着祂的门徒所行的征兆，却使他们充满困扰。因为那些被认为偷走尸体的门徒，被发现使别人的尸体复活。但恶人惊慌地说——\Quote{祂的门徒偷走了祂的身体；}这样当真相被发现时，他们便会遭人轻视。但门徒们（他们[据称]从活着的卫兵那里偷走了尸首），却被发现奉被偷者的名攻击死亡；好让[死亡]不再偷走活人的生命。因此，在祂被钉十字架之前，祂使聋子听见，为的是在祂被钉十字架之后，众人都能耳闻并相信祂的复活。因为祂事先借着哑巴的开口（其口被开启）坚定了我们的听觉，使它对圣言的宣讲不再怀疑。我们的救赎者样样装备齐全，为能样样将我们从俘虏者手中解救出来。因为我们的主岂止披戴了身体，还穿上了肢体和衣裳；好借着祂的肢体和衣裳，使患病者得着鼓励去接近医治的宝库，使被祂仁慈鼓励的人接近祂的身体，而被祂威严吓倒的人接近祂的衣裳。因为有一个妇人，她的恐惧只容许她接近祂的衣边；\BibleRef{玛 9:20} 但另一个妇人，她的爱却驱使她接近祂的肉身。借着那因祂衣裳得医治的妇人，那些未因祂言语得医治的人被羞愧；而借着那亲吻祂脚的妇人，那不愿亲吻祂嘴唇的人被斥责。\par

14. 我们的主借着微小的事物赐下伟大的恩赐；为教导我们那些轻视大事的人被剥夺了什么。因为若从祂的衣边就能如此偷偷窃取医治，难道祂以明确的话语赐予医治时不能医治吗？若污秽的双唇因亲吻祂的脚而成圣，洁净的双唇岂不更应因亲吻祂的口而成圣？因为那有罪的女人借着亲吻获得了祂圣足的恩宠，这双圣足带着辛劳前来，为带给她罪的赦免。她白白地用油膏抹医治者的脚，因为祂白白地为她的病带来了医治的宝藏。因为那满足饥饿者的一位作客，并非为了自己的肚子；而是为了罪妇的悔改，那成义罪人的一位使自己成了客人。\par

15. 因为我们的主所渴求的并非法利塞人的美味，而是罪妇的眼泪。因为当祂因所渴求的眼泪而满足并苏醒时，祂转身斥责那请祂赴宴、宴席会朽坏的人，为表明祂作客并非为了身体的食物，而是为了灵魂的帮助。因为我们的主与饕餮之徒和酒徒为伍，并非如法利塞人所想的那样为了享乐；而是在他们作为必朽者的饮食中，为祂们混入祂的教导作为生命之药。因为正如在饮食的事上，恶者向亚当及其配偶献上了致命的计谋，同样在饮食的事上，善主向亚当的子孙献上了赐生命的计谋。因为祂是下来捕捞失丧生命的渔夫。祂看见税吏和娼妓纵情于挥霍和酗酒；祂急忙在他们的聚集之处撒下网，为要从使身体肥胖的食物中捕获他们，转向使灵魂健壮的禁食。\par

16. 法利塞人为我们的主在宴席上做了盛大预备；而罪妇在那里为祂做的事却很少。然而，他借着他的美味显出了他对我们的主爱的渺小；而她借着她的眼泪显示了她对我们的主爱的伟大。因此，那邀请祂赴盛大宴席的人，因爱的渺小而受斥责；而她借寥寥数滴眼泪赎偿了她诸多过犯的愚妄。西满法利塞人将我们的主看作先知——是由于征兆，而非由于信德。因为他是以色列之子，当征兆临近时，他也临近征兆之主；当征兆止息时，他也赤身无信地站立。此人见我们的主行征兆，便尊祂为先知；但我们的主止息征兆时，他族人的疑惑之心便进入了他。\Emph{这人若是先知，必知道这女人是个罪妇。}但我们的主，在每一处一切都容易的一位，在这里也并未止息祂的征兆。因为祂看见，因祂稍止征兆，法利塞人盲目的心就背离了祂。因为他错谬地说：\Emph{这人，若是先知，必知道}。因此，在此思虑中，法利塞人对我们的主疑惑不定，不知祂是否是先知；但正是借此思虑，他得知了祂是先知之主；如此，从错误进入他的源头，我们的主从那里向他伸出援手。\par

17. 主随即向他说了两个债务人的比喻，并让他作审判官，为要藉他的口捉住那心里没有真理的人。\Emph{一人欠五百个德纳}。于是主向法利塞人显示那有罪妇人的诸多过犯。那曾猜想主不知道她是罪人的人，结果从主口中听见她的罪债何其大。法利塞人曾猜想主不知道那妇人是谁，也不知道她有罪的名声，结果发现自己竟不知道主是谁，也不知道主的名声。如此，他在自己的错误中受责备，他甚至未曾察觉自己的错误。因为他确知自己在错误中的知识，在错误中逃避了他。但他从那位来提醒犯错者的主那里得到了提醒。法利塞人见过主所行的大奇迹，正如以色列见过梅瑟所行的；但因他心中没有信德，他所见的那些奇事未能与信德结合，一点小事就阻碍并取消了它们。\Quote{这人若是先知，必知道这女人是个罪人}。因为他忽略了他所见的奇事，盲目轻易进入了他。他本是以色列子民，有可畏的征兆伴随他们直到海边，使他们敬畏；在旷野中有福的奇迹环绕他们，使他们与天主和好；但因缺乏信德，为了一点原因，他们就弃绝了这些事，说：\Quote{至于这个领我们上来的梅瑟，我们不知道他遭遇了什么事}。\BibleRef{出 32:1} 因为他们不再留意环绕他们的伟大作为。他们察觉梅瑟不在他们近旁，以致于因这个临近的原因，他们就趋向埃及的异教。梅瑟暂时从他们面前被移开，好使那在他们面前的牛犊显现，使他们也能公开朝拜它；因为他们早已在心中秘密朝拜它了。\par

18. 但当他们的异教从内在变为公开时，梅瑟也由隐藏而公开显现，为要公开惩罚那些在荫庇他们的圣云之下欢乐地行异教的人。但天主将羊群的牧者从他们那里撤去四十天，好让羊群显明它的信赖是寄托在牛犊上。当天主以各种欢乐喂养羊群时，羊群却为自己选择了那连吃都不能的牛犊作为牧者。那令他们敬畏的梅瑟被移走，好使偶像崇拜在他们口中高声呼喊，而梅瑟的约束曾将它压制在他们心中。因为他们呼喊：\Quote{为我们造神，在我们前面引路}。\BibleRef{出 32:1}\par

19. 但当梅瑟下来，看见他们在广阔的平原上以鼓和铙钹欢乐地行异教。他立即藉肋未人和出鞘的刀剑使他们羞愧。同样，在此处，主暂时隐藏祂的知识，当那有罪妇人走近祂时，好让法利塞人形成他的思想，正如他的祖先形成了那邪恶的牛犊。但当法利塞人的错误在他内发展到顶点时，主的知识便显现出来，驱散了它：\Quote{我进了你的家；你没有给我水洗脚：但她却用眼泪湿润了我的脚。所以，她的罪，许多的罪，都赦免了}。\BibleRef{路 7:44} 但法利塞人听见主提到那妇人的罪，\Emph{许多的罪}，便大大羞愧，因为他大大错了。他曾以为主甚至不知道她是罪人。主先前曾显出好像不知道她是罪人的样子。因为祂允许那见过祂记号的人，显露他心中的疑惑，好显明他的心被束缚于他祖先的不敬虔之中。但医师，藉他的药物引出隐藏的疾病，不是疾病的帮助者，而是它的毁灭者。因为当疾病隐藏时，它在肢体中掌权；但当被药物显明时，它就被根除。同样，法利塞人见了大事，却在小事上怀疑。但当主看见他的渺小使他在心中轻视大事，祂立刻向他显明她不但是罪人，而且她的罪甚多，好让他因小事而羞愧——他这不信奇事的人。\par

20. 天主给了以色列空间，让它在广阔的旷野中扩大其异教；天主用磨快的刀剑将其剪除，免得他们的偶像崇拜在外邦人中传开。同样，主允许法利塞人想象乖谬的事，为要反过来适时责备他的骄傲。因为关于那有罪妇人做得对的事，法利塞人想得错了。但主反过来责备他，关于他错误地扣留的正当之事：\Quote{我进了你的家；你没有给我水洗脚}。看哪，扣留了应尽的本分！\Quote{但她用眼泪湿润了它们}。看哪，偿还了应尽的本分！\Quote{你没有用油抹我}。看哪，疏忽的记号！\Quote{但她用香膏抹了我的脚}。看哪，热忱的记号！\Quote{你没有吻我}。看哪，敌意的见证！\Quote{但她没有停止吻我的脚}。看哪，爱的记号！如此，主藉这列举表明，法利塞人欠主所有这一切，却扣留了它们；但那个有罪妇人进来，偿还了他所扣留的一切。因为她偿还了那错误扣留之人的债务，义者便赦免了她自己的债务，即她的罪。\par

21. 法利塞人怀疑主不是先知时，无意中为真理作了见证，说：\Quote{这个人若是先知，必知道这女人是罪人。}因此，如果查出主知道她是罪人，那么按你的话，法利塞人啊，祂就是先知。于是主急忙表明她既是罪人，她的罪又很多，好让他亲口说的话反驳他自己是个说谎者。因为他是那些说\BibleQuote{除了天主一个外，谁能赦罪呢？}（见：\BibleRef{谷 2:7}）的人的同伙。主从他们那里得到见证，因此，那能赦罪的就是天主。从此，争论点就是：主应当向他们显明祂是否能赦罪。所以祂迅速医治了可见的肢体，好证实祂赦免了不可见的罪。因为主把那些期待抓住发言者的话摆在他们面前，当他们照自己的意愿冲上前来抓住祂时，他们反而照祂的意愿被祂抓住。\BibleQuote{放心，孩子！你的罪赦了。}（见：\BibleRef{玛 9:2}）当他们急匆匆要以亵渎的罪名捉拿祂时，他们无意中为真理作了见证。因为\Quote{除了天主一个外，谁能赦罪呢？}于是主驳斥他们【好像说】：\Quote{如果我显明我能赦罪，即使你们不相信我是天主，但你们仍要信守你们的话，即断定凡赦罪的就是天主。}因此，为教训他们祂赦免罪，祂赦免了那人隐藏的罪，并使他公开地拿起床来行走；好叫那抬着〔躺卧之人〕的床被抬起来，他们就会相信那致死的罪已被除去。\par

这是一件奇妙的事：主在那里自称人子，祂的敌人们却无意中因祂赦罪而承认祂是天主。因此，当他们自以为用诡计网罗了祂时，祂反而用他们的诡计缠住了他们；祂使这成为祂真理的见证。所以，他们的恶念成了他们沉重的桎梏；为使他们无法摆脱自己的桎梏，主就赐力量给祂对他说话的那个人——\BibleQuote{起来，拿起你的床，回家去吧！}（见：\BibleRef{玛 9:6}）因为这见证不能再被推翻，好像祂不是天主；因为祂赦免了罪。也不能虚假地断言祂没有赦罪；因为看！祂医治了〔人的〕肢体。因为主把祂隐藏的见证与显明的见证联结起来，好叫他们自己的见证窒息不信的人。因此，主使他们的思想攻击他们，因为他们曾与善者争战，祂却以医治的能力与他们的疾病争战。因为法利塞人西满所想象的，以及他的同伴经师们所想象的，他们在心中秘密地想象；但主公开地陈述出来。主将他们隐藏的想象摆在他们面前，好让他们知道，祂的知识启示并显示他们的隐秘事；这样，即使他们未因祂公开的标记认识祂，但当祂揭示他们隐秘的想象时，他们或许能认识祂；并且，即使只凭这一点——祂洞察了他们的心——他们的心也能察觉祂是天主；至少当他们看见他们的想象不能向祂隐藏时，他们可以停止对祂怀有恶念。因为他们心中怀有恶念；但祂公开揭露出来，藉着这句话：\Quote{你们为什么心怀恶念呢？}这样，因主觉察他们隐秘的想象，他们应该认识祂隐秘的天主性。因为那天主性，正是因他们在错误中亵渎它，才藉着那亵渎得以向他们显明。因为他们在肉身中亵渎主，以为祂不是天主，将祂从高处推下低处；但藉着那肉身，祂被他们认识为天主，藉着那在他们中间来往的身体。因为他们将祂推下深渊，企图证明这一点：在上面的天主，不能以肉身方式生在下面。但祂藉着祂上升到高处，教导他们这一点：即使是被送到下面的身体，它的本性也不是上升高处而非下降深处；好藉着那从下面升到高空中的身体，他们可以认识天主：因祂的恩宠，祂从高处降到了低处。\par

22. 但为何我们的主不以严厉的斥责，反倒向那法利塞人讲了一个劝说的比喻呢？祂温柔地对他讲这个比喻，为使他——虽然乖戾——不知不觉地被诱而纠正自己的乖张。因为被冷风之力冻结的水，太阳的温热轻轻融化。同样，我们的主没有立刻严厉地反对他，恐怕给悖逆者以再次悖逆的机会。反而用和蔼的方式使他受轭，以便在受轭之后，即使他悖逆，也能按祂的旨意与他一同工作。如今，因为西满心存骄傲，我们的主便以谦卑开始与他交谈，为使自己不照他的愚妄成为他的老师。因为如果那法利塞人持守法利塞人的骄傲，我们的主怎能使他获得谦卑，当谦卑的宝库不在他手中时？但既然我们的主教导众人谦卑，祂便表明祂的宝库远离一切形式的骄傲。这是为了我们的缘故，为教导我们：无论骄傲进入什么宝库，它都是通过自夸进入的。为此，\BibleQuote{不要让左手知道你右手所做的。}\BibleRef{玛 2:3} 我们的主那时没有采用严厉的斥责，因为祂的来临是出于恩宠；祂并未完全免去斥责，因为祂后来的来临将是出于报应。因为祂在谦卑的来临时就使人畏惧，因为\BibleQuote{落在祂手中是多么可怕的事}\BibleRef{希 10:31}，当祂\BibleQuote{在火焰中}\BibleRef{得后 1:7}降临时。但我们的主多半是以劝说的方式而非斥责来施予帮助。因为温和的雨水滋润大地并渗透其全部；而暴雨则紧压并硬化地面，使地面无法吸收。因为严厉的言语激起愤怒，而随之而来的是冤屈。当严厉的言语打开门时，愤怒便进入，紧跟着愤怒，冤屈也一同进入。\par

23. 但由于所有帮助都伴随着谦卑的言语，那位来施行帮助的便采用了它。请观察谦卑之言的能力何等强大；因为看！激烈的愤怒借此被平息，膨胀心灵的波涛借此被平息。但请听这话从何而来。那法利塞人想：\BibleQuote{这人若是先知，必知道}。这里可见轻蔑和亵渎。请听我们的主在答复中如何应对：\BibleQuote{西满，我有句话对你说}。这里可见爱和斥责。因为这是朋友间所用的爱的话语。因为当仇敌责备仇敌时，他不会像这样对他说话；因为愤怒的疯狂不允许仇敌互相理智地说话。但那为钉死祂的人祈祷的那一位，为表明愤怒的狂怒对祂毫无能力，将要审问那些钉死祂的人，以表明祂受理性而非愤怒的支配。\par

24. 因此，我们的主在谈话开始时放了一句和解的话，为通过和解使那法利塞人平静，在他的心中已进入不和与分裂。祂是医生，针对损害中的人安排祂的疗法。我们的主便射出这句话如箭，并以和解为箭镞装在箭头上。祂又以爱来涂抹它，这爱使肢体安宁；以致当它飞入那充满不和之人的心中时，他立刻从不和转为和谐。因为一听到我们的主谦卑的声音，说——\BibleQuote{西满，我有句话对你说}，那暗中的藐视者便回答：\BibleQuote{请说，主}。因为甘甜的声音进入他苦涩的心，并从其中生出愉快的果实。因为在这声音之前他是暗中藐视者，在这声音之后他成了公开尊敬者。因为谦卑以其甘甜的言词，甚至使仇敌屈服而尊敬它。因为谦卑不是对朋友测试它的能力，而是在仇敌身上展示它的胜利。\par

25. 如此，天上的君王以谦卑的铠甲武装自己，于是征服了那苦涩者，并从他那里得了一个好的答复，作为确凿的胜利保证。这就是保禄所说的铠甲，通过它\BibleQuote{我们攻破那高傲地起来反对认识天主的堡垒}\BibleRef{格后 10:5}。因为保禄在自己身上领受了这铠甲的证明。因为正如他曾在骄傲中作战，却在谦卑中被征服，同样，每一个\Quote{自高自大起来}反对这谦卑的高傲之物都将被征服。因为扫禄曾去用严厉的言语压迫门徒，但门徒的老师却用一句谦卑的话征服了他。因为当那位无所不能者向他显现时，祂放弃其他一切，只以谦卑对他说话，为教导我们：柔软的舌头比一切更能对付刚硬的心思。因为保禄没有听到威胁或恐吓的话，而是听到软弱无力、不能报复的话：\BibleQuote{扫禄，扫禄，你为什么迫害我？}\BibleRef{宗 9:4} 但那些被认为甚至不能报复的话，却藉着将他从犹太人那里拉出来，使他成为一件美器，而施行了报复。他那充满犹太人苦涩意志的人，随后被十字架的甘甜宣讲所充满。当他充满钉死者的苦涩时，他在他的苦涩中摧残教会。但当他充满被钉者的甘甜时，他使钉死者的会堂感到苦涩。我们的主于是以谦卑的声音与那曾用强硬锁链攻打祂教会的人争战。如此，扫禄——他曾用苦涩的锁链捆绑门徒——被愉快的劝诱所捆绑，免得他再将门徒投入锁链，因为他已被被钉者所捆绑，祂止息了恶言，所有对抗祂的人都不能捆绑或伤害祂。但当保禄停止捆绑门徒时，他自己却被迫害者用锁链捆绑。但当他被锁链捆绑时，他藉着自己的锁链解开了偶像崇拜的锁链。\par

26. \Emph{扫禄，扫禄，你为什么迫害我？}那位在下界已征服了祂的迫害者，并在上界统治天使的，从高处用谦卑的声音发言。祂在地上时曾向钉死祂的人宣告了十个祸哉，但当祂在天上时，却没有向祂的迫害者扫禄宣告一个祸哉。现在，我们的主向钉死祂的人宣告祸哉，是为了教导祂的门徒不要因杀害祂的人而惊惶。但我们的主从天上谦卑地发言，为的是使祂的教会的首领们也能谦卑地发言。如果有人要说：\Quote{我们的主在什么事上谦卑地与保禄说话呢？因为，看！保禄的眼睛严重受损了；}让他知道，这惩罚并非出自我们仁慈的主——祂谦卑地说了那些话——而是出自那强烈照射的强烈光芒。并且这光照击保禄，并非因他的行为作为报应，而是由于它的强光伤了他，正如他也说过：\Emph{我站起来，什么也看不见，因为那光太荣耀了。}\BibleRef{宗 22:11}但如果那光是荣耀的，保禄啊，荣耀的光怎么反而成了使你目盲的光呢？那光是按其本性照亮上界的，但反其本性而在下界照耀。当它照明上界时，它是愉悦的；但当它照耀下界时，却是使人目盲的。因为那光既是痛苦的，又是愉快的。它对肉体的眼睛是痛苦而猛烈的；对属于火与灵的人却是愉快而光明的。\BibleRef{玛 4:11}\par

27. \Emph{因为我看见一道从天上来的光，比太阳更亮，它的光照耀着我。}\BibleRef{宗 26:13}于是，强烈的光线无节制地倾泻出来，涌向软弱的眼睛，而适度的光线本是滋养眼睛的。因为，看！太阳也适度地助益眼睛，但过度和无节制则会伤害眼睛。并且它击打眼睛，不是出于愤怒中的报复。因为，看！它是眼睛的朋友，为眼球所喜爱。这是一件奇妙的事：它以其柔和的光泽与眼睛为友并助益眼睛，却以其强烈的光线敌视并伤害眼球。但如果这在下界的太阳，与在下界的眼睛有相似的本性，尚且伤害它们——出于强烈而非愤怒，出于其自然力而非愤怒——那么，那从上来的光，与上界事物相似，岂不更以其强烈伤害一个在下界的人，他突然注视了与他本性不相似的事物呢？因为，既然保禄可能因他所习惯的太阳的强烈而受伤——如果他不是按照习惯注视它——那么，他岂不更因那光——他的眼睛从未习惯——的荣耀而受伤吗？因为，看，达尼尔也\BibleRef{达 10:5}被熔化，向四面倾流，在天使的荣耀面前，天使的强烈光芒突然照在他身上！并且，他的软弱人性被熔化，并非因为天使的愤怒，正如蜡烛在火前熔化，并非因为火的愤怒或敌意，而是因为蜡烛的软弱，不能在火面前保持坚固站立。因此，当两者靠近时，火的力量因其性质而占优势；而蜡烛的软弱则被降低，甚至低于它原有的软弱。\par

28. 但天使的威严在其自身彰显；血肉的软弱在其自身无法承受。因为\Emph{我的内脏变为腐朽。}\BibleRef{达 10:8}但是，人们看见他们的同类，就在他们面前昏厥：然而，他们并非因对方的光辉夺目而动摇，而是因对方严厉的意志。因为仆役被主人的愤怒惊吓，受审判的人因对审判官的恐惧而战栗。但达尼尔遭遇这事，并非由于天使的威胁或愤怒，而是由于他可畏的本性和压倒性的光明。因为天使到他那里去，不是带着威胁。因为如果他是带着威胁而来，一张满是威胁的口怎能变为充满平安，当它来说：\Emph{愿你平安，你}蒙爱的人？因此，那曾是雷霆之源的口——\Emph{因为他言语的声音好像大军的声音}，\BibleRef{达 10:6}——那声音对达尼尔成了充满并涵容平安的源泉。当那声音到达那惊恐的耳朵——它们渴望那鼓励性的平安问候——时，便有平安的饮品倾流而出（为他）。并且，藉天使后来的平安（话语），那先前被他先前的声音所惊吓的耳朵受到了鼓励。因为（他说）：\Emph{愿我主发言，因为我已坚强起来。}\BibleRef{达 10:19}但因在那感人至深的异象中，那火炽的天使即将宣告关于祂（主）的什么事也未曾有，为此，那威严（天使）便主动向那卑微（先知）致以平安的问候；好藉那可畏的威严所给予的喜乐问候，除去那盘踞在卑微心灵上的恐惧——那心灵是惊恐的。\par

29. 但关于那天使的主，我们该说什么呢？祂曾对梅瑟说：\Quote{\Emph{没有人看见我还能存活}}。\BibleRef{出 33:20}是因祂怒气的猛烈，凡看见祂的人就必死吗？还是因祂存有本体光辉？那本体并非受造，也非被造，以致受造和被造的眼睛不能瞻仰。若因祂的怒气，凡看见祂的人不能存活，看！祂因对梅瑟的大爱，本会准许他看见祂。因此，自存者\OriginalTerm{自存者}{Self-Existent}的显现杀死那些观看祂的人；但祂并非因严酷的怒气，而是因祂强大的光辉而杀。为此，祂因大爱准许梅瑟看见祂的荣耀；又同样因大爱阻止他看见祂的荣耀。但这并非祂尊威的荣耀有丝毫减少，而是软弱的眼睛不足以承受祂荣耀的汹涌波涛。因此，天主出于爱，愿意梅瑟的视线注视祂荣耀的美妙光辉，也出于爱，不愿意梅瑟的视线在祂荣耀的强力光芒中被蒙蔽。所以梅瑟看见了，又没有看见。他看见了，为能被提升；他没有看见，为能不受伤害。因他所看见的，他的卑微被高举；因他所没有看见的，他的软弱未被蒙蔽。正如我们的眼睛看太阳，却又没有看它；因所见而得益，因未见而不受损。同样，眼睛观看，为能获益；但它不敢直视，为免受伤。所以，天主因爱阻止梅瑟看见那过于他眼睛所能承受的荣耀；正如梅瑟因爱阻止他本族人民看见那过于他们眼睛所能承受的光辉。因为他从那遮盖他、伸手遮蔽他、隐藏光芒以免伤他的那一位那里学到，他也该蒙上帕子，向软弱者隐藏那过于强烈的光辉，以免伤害他们。当梅瑟看到有死之肉的子孙不能注视他脸上那借来的荣耀时，他心中沮丧；因为他曾试图大胆瞻仰永恒存有者的荣耀；在那荣耀的洪流中，看！天上和地下都沉浸其中并涌出；其深无人能测，其岸无人能及，其终无极尽。\par

30. 若有人说：\Quote{难道天主不可能使梅瑟看见那荣耀而不受伤，同样也使保禄看见那光而不受害吗？}说这话的人要明白：虽然天主的权能和驾驭之力能改变眼睛的本性，但混淆自然的秩序与天主的智慧和本性不符。因为，看！工匠的手也能轻易毁坏他的作品，但破坏美好的装饰与工匠的明智不符。若有人想谈论某件在他看来似乎合宜的事，说：\Quote{天主做这事是合宜的，}就要知道：他自己不说这样关于天主的话才是合宜的。因为一切合宜事之首便是：人不应教导天主什么是合宜的。因为人成为天主的导师是不相宜的。这是极大的邪恶：我们竟成为那位的导师，而我们这些受造的口舌连祂手工作品的形成也无法述说。因为口舌凭胆量教导那位天主什么是合宜的，这是不可宽恕的罪过——这口舌原是靠祂的恩宠才学会说话。若有人说：\Quote{天主做这事本是合宜的，}我因有口有舌，也可以说：\Quote{天主本不应赐给人自由，使人借此辱骂那不可辱骂者。}但我不敢说祂赐予自由是不合宜的，免得我也成为那不可受教者的导师。因为祂是正义的，若不赐给人自由，祂自己就会受责难，仿佛出于吝啬而向卑微的人扣留了那使人伟大的恩赐。因此祂及时赐予自由，为使自己不被正义地责难；尽管藉着自由——祂自己的恩赐——看！亵渎者邪恶地辱骂祂。\par

31. 为何梅瑟的眼睛因他所见的荣耀而发光，相反，保禄的眼睛非但没有发光，反而完全失明？然而我们可以肯定，梅瑟的眼睛并不比保禄的更强，因为它们在血与肉的同一兄弟情谊中是相似的。但另一种能力藉恩宠支撑了梅瑟的眼睛；而保禄的眼睛，除了其自然能力之外，仁慈地没有加添任何能力，反倒愤怒地剥夺了它们。但若我们说他们的自然能力被剥夺了，正因如此他被打败并被那超越的光所征服——因为若他们的自然能力保留，他们本可忍受那超自然的光。然而让我们确信：每当有超然之事显现，超越并超过我们的本性，我们的自然能力就不能在其面前站立。但若另一方面，一种超越我们自然能力的能力加给我们，那么藉着这超乎本性额外领受的能力，我们就能在超自然地临于我们身上的任何奇异事物面前站立。\par

32. 请看！我们耳目之能力在我们内，以自然方式形成；然而我们的视觉与听觉在强大的雷声和闪电前无法站立；首先，因为它们带着猛烈而来；其次，因为它们的威力突然使我们软弱的人惊讶而惊愕。这正是保禄所遇到的。因为光芒的威力突然使他的衰弱眼睛惊讶并伤害它们。但声音的伟大降低他的力量，进入他的耳朵并打开了它们。因为他们已被犹太人的争辩如蜡一般封住。因为声音并未犁开耳朵，如同光芒伤害眼球。为什么？而是因为适合他听见，却不应看见。因此听觉之门由声音如钥匙一般被打开；但视觉之门却被应该打开它们的光芒关闭。那么为什么适合他听见？显然因为藉着那声音，我们的主能显示自己正受撒乌耳的迫害。因为祂不能以视觉显示自己正受迫害；因为没有途径可使达味之子被看见逃跑而撒乌耳追赶祂。因为这事确曾发生在第一位撒乌耳和第一位达味身上。一个在追赶，另一个在被迫害；他们彼此看见并被看见。但在此处，唯独耳朵能听到达味之子所受的迫害；眼睛不能看见祂受迫害。因为祂是以他人的身份受迫害，而祂自己却在天堂——祂先前在地上时曾亲身受迫害。因此（撒乌耳的）耳朵被打开，眼睛被关闭。祂虽然不能以视觉在撒乌耳面前显示自己受迫害，却以言语在他面前显示自己受迫害；当时祂喊道：\BibleQuote{撒乌耳，撒乌耳，你为什么迫害我？}因此，他的眼睛被关闭，因为它们不能看见基督的迫害；但他的耳朵被打开，因为它们能听见祂的迫害。同样，梅瑟的眼睛虽是属肉体的眼睛，如同保禄的，但他的内在眼睛却是基督徒的；因为\BibleQuote{梅瑟写论及我}\BibleRef{若 5:46}，但保禄的外在眼睛是睁开的，而内在（眼睛）是关闭的。那时因为梅瑟的内在眼睛明亮照耀，他的外在眼睛也被造得明亮照耀。但保禄的外在眼睛被关闭，为叫藉着外在的关闭，能实现内在的开启。因为那藉外在眼睛不能在祂的标记中看见主的人，当那些属肉体的眼睛关闭时，他藉内在的眼睛看见了。又因为他亲身接受了证据，他写信给那些肉体眼睛充满光明的人——\BibleQuote{愿祂照亮你们心目的眼睛}\BibleRef{弗 1:18}。因此，显于犹太人外在眼睛的标记，对他们毫无益处；但心中的信德打开了外邦人心中的眼睛。但因为，若梅瑟以平常形貌从山上下来，没有那面容的光辉，却说：\Quote{我在那里看见了天主的荣耀，}那些无信德的祖先不会相信他；同样，若保禄没有遭受眼睛的失明，却说：\Quote{我听见了基督的声音，}那些钉死基督的子孙不会接受它为真。因此，祂在梅瑟身上，出于慈爱，安置了卓越的光辉标记，好叫欺骗者相信他看见了天主的荣耀；但在撒乌耳身上，如同对迫害者，祂安置了可憎的失明标记，好叫说谎者相信他听见了基督的话；好叫你不再反对梅瑟，也使他们不怀疑保禄。因为天主在瞎子的身体上安置标记，并派遣他们到那些在错误中、习惯在衣边上作标记的人那里。但他们不记得衣服上的标记，而在身体的标记上大大犯错。那些看见梅瑟荣耀的祖先，并没有听从梅瑟；那些看见保禄失明的子孙也没有相信保禄。但在旷野中，他们三次威胁要用石头掷死梅瑟和他的家，如同对待狗。因为\BibleQuote{全会众吩咐用石头掷死他们}\BibleRef{户 14:10}。他们三次用棍棒鞭打保禄，如同狗一般打他的身体。\BibleQuote{三次我受了棍打}\BibleRef{格后 11:25}。这些就是因爱他们的主而如狗般被打、如羊群般被撕裂的狮子；那些曾用石头掷死看护他们的牧者的羊群，好叫贪食的狼统治他们。\par

33. 但是那些用贿赂收买士兵的钉死主的人，或许曾就保禄说过：\Quote{门徒们用贿赂收买了他；因此他与门徒们来往。} 因为那些藉行贿而力阻主复活被宣讲的人，以贿赂之名毁谤保禄，使他的启示不被相信。因此，那声音使他惊惶，那光使他失明，好叫他因惊惶而平息其暴烈，他的失明则使毁谤者蒙羞。因为那声音如此温良地对他说：\Emph{扫禄，你为什么迫害我？}；这使他听觉得到惊惶；那光使他视觉失明，为的是当毁谤者说他收了贿赂，因而被收买说谎时，他因那光而导致的失明能反驳他们，表明正是藉这光，他才被迫说出真理。这样，那些以为他的手受了贿赂、因而口舌说谎的人，就能知道他的眼已失去光明，而因这光明，他的口舌反而宣告了真理。但又有另一个原因，温良的声音伴随着压倒一切的光：即好像从温良到达高举，主为迫害者产生了帮助；正如祂所有帮助的产生，都是从卑下到伟大。因为主的温良从母胎延续至坟墓。并且请注意，伟大紧随着祂的卑下，高举紧随着祂的温良。因为当祂的伟大在各方面被观察到时，祂的天主性藉荣耀的标记被启示出来，好使人知道那站在他们中间的，不是一位，而是两位。因为祂的性体不是单单卑下的性体，也不是单单高举的性体；而是有两个性体相合，一个与另一个：即高举的与卑下的。\par

因此，这两个性体显出它们的特质；好使人类藉二者各自的品质，能分辨二者；以免人以为祂仅仅是——一位因相合而为二者；但要使人知道祂在混合方面是二，尽管在存有方面为一。我们主藉祂的谦卑与高举，也在往\Place{大马士革}的路上教导了\Person{保禄}这些事。\par

34. 因为主以温良向扫禄显现，因为温良接近祂的伟大；好因祂的伟大使人知道那温良说话的是谁。因为正如祂的门徒在地上宣讲我们的主，既在温良中也在高举中——在祂受迫害的温良中，也在祂标记的高举中——同样，主也在保禄面前以温良和以高举宣讲了祂自己——在闪光的强大能力的高举中，也在那温良声音的温良中，那声音说：\Emph{扫禄，你为什么迫害我？}——好使祂的门徒在众人面前所宣讲关于祂的宣讲，应与祂自己所宣讲关于自己的宣讲相似。但正如，如果祂没有温良地说话，那里就不会知道祂是温良的；同样，如果祂没有以压倒性的光显现，那里就不会知道祂是高举的。\par

35. 如果你说：\Quote{有什么必要祂要谦卑地说话？祂岂不能也藉光的大能说服他吗？} 你要知道，你这发问的人，这个答辩可以还给你：因为祂必须谦卑地说话，所以祂便谦卑地说话。因为那在一切事上有智慧者，在那里没有行任何不当行的事。因为那赐给工匠知识，叫他们用适合的器皿分别作各事者，祂自己岂不知道祂所赐给别人认识能力的事吗？所以，凡由天主性所行或正行的事，正是祂在那时所行的事，是为了促进[天主]在那时的工作，尽管对盲目者而言，天主的安排似乎相反。但为了避免我们用言语的约束来限制一个明智的探问者，一个愿意像种子靠雨滴生长那样因真正的说服而成长的人，请你知晓，探问者啊，因为扫禄是迫害者，但我们的主竭力使他成为受迫害者而非迫害者，所以祂因自己的智慧急忙呼喊——\Emph{扫禄，你为什么迫害我？}——为的是，当扫禄正被造就为门徒时，听到那造就他为主的说：\Emph{你为什么迫害我？} 他就能知道那他将成为其仆人的主人，是一位受迫害的主人，因而能迅速摆脱他先前主人们的迫害，并穿上他受迫害的主人的受迫害状态。如今，任何希望教人某事的主人，要么藉行为要么藉言语教导他。但如果他既不藉言语也不藉行为教导他，那人就不能受教于他的技艺。所以，尽管我们的主没有藉行为教导保禄谦卑，但却藉声音教导了他忍受迫害，这是祂不能藉行为教导的。因为我们的主被钉十字架之前，祂藉行为教导门徒谦卑地忍受迫害。但祂藉十字架完成了迫害之后，正如祂所说：\Emph{看！一切都完成了。}\BibleRef{若 19:30} 祂不能徒然返回，重新开始那一次永远明智地完成的事。或者你为什么再寻求天主之子的十字架与羞辱呢？\par

36. 因为尽管我们的主在祂的恩宠中，先前已将祂天主性的威严降入卑微，但日后在祂的公义中，祂不愿再将那已被提升的人性的渺小重新带回屈辱。然而，由于那迫害的门徒必须学会忍受迫害，而师傅却不可能再次降下并重新受迫害，祂便用言语教导他那些无法以行动教导的事。\BibleQuote{扫禄，你为什么迫害我？} 这句话的解释是这样的——\Quote{扫禄，你为什么不在我内受迫害？} 但为使扫禄不致以为我们的主是因软弱而受迫害，那照在他身上的压倒性光芒的力量便说服了他。因为如果扫禄的眼睛不能忍受那光芒的照耀，扫禄的手又怎能捆绑和锁住那光芒之主的门徒呢？他曾用双手捆绑门徒，好在他被捆绑时学到他们的力量；而他的眼睛却不能忍受光芒，好让他在光芒的力量中认识自己的软弱。然而，当主对他说：\BibleQuote{扫禄，你为什么迫害我？}时，那光芒的力量岂非曾照在他身上？那时，鉴于保禄当时所怀骄傲的疯狂，他或许会对祂说：\Quote{我如此迫害你，是因为你说了：\InnerQuote{你为什么迫害我？}因为谁会不迫害你呢？你竟以如此力量，用这些微弱的喊声扰乱你的迫害者。} 但主的谦卑在声音中被听见，光芒的力量在光束中显现。因此，保禄不能因着光芒的荣耀而轻视那声音的谦卑。\par

37. 于是，他的耳朵因所听见的声音而成为门徒，因为他的眼睛不足以承受他们所看见的光芒。那光芒破晓的神奇景象照射在他的眼球上，伤害了它们；而光芒之主的言语进入他的耳朵，却未造成伤害。但光芒与其主之间，哪一样应当更强大呢？因为如果那由祂受造的光芒如此压倒一切，那么创造这光芒的祂，又该是何等更加强大呢！但如果光芒之主是压倒一切的——祂确实如此——那么祂的声音如何能进入听觉而不伤害它，如同那光芒伤害了视力一样呢？但请听我们的主因祂的恩宠所行的奇事与神迹。因为我们的主不愿贬低属于祂的光芒；但祂，作为光芒之主，反倒谦卑了自己。然而，正如光芒之主大于属于祂的光芒，同样，光芒之主谦卑自己而非贬低光芒，是何等大的荣耀。\par

38. 同样，在夜晚，当祂祈祷时，经上记着：\BibleQuote{有一位天使从天上显现给祂，加强祂的力量。}\BibleRef{路 22:43} 但在这里，天上和地上的所有口唇，都不足以感谢那创造天使的祂；因为祂为罪人的缘故，竟由祂亲手所造的天使所加强。正如那天使从高天荣耀光辉地站立，而天使之主为了使被贬低的人得以\Emph{升高}，却站在卑微和屈辱中；同样，此处那光芒发显显耀；但光芒之主为帮助一个迫害者，却以谦卑的声音和卑微的话语说话。\par

39. 因此，那压倒一切的光芒，因未被减弱，便以压倒性的显耀进入眼球，伤害了它们。但光芒之主，因祂为了帮助而降低了自己，祂卑微的声音便进入那有需要的耳朵，帮助了它们。然而，为使那已变得卑微的声音的帮助不致失效，那光芒的力量并未被降低；这样，因着那未被降低的光芒，那被降低的声音的帮助才得以被相信。但这是一个奇迹：直到我们的主在声音中变得卑微之前，保禄在行为上并未变得卑微；因为正如我们的主在降生成人、披上身体之前，与祂的圣父同在高天；但在祂的高举中，人们并未学到谦卑；而当祂谦卑自己，从高天降下时，借着祂的谦卑，谦卑很快便在人中传播；同样，在祂复活升天之后，祂在天主父的右边享受荣耀，但借着那高举，保禄并未学到谦卑。因此，那被高举、坐在圣父右边的祂，停止了荣耀崇高的言语，如同受冤屈、受压迫者一般，以微弱温柔的话语呼喊说——\BibleQuote{扫禄，扫禄，你为什么迫害我？} 这样，谦卑的话胜过了严酷的缰绳。因为借着谦卑的话，如同借着缰绳，那被迫害者将迫害者从迫害者的宽路领入了被迫害者的窄路。既然所有奉我们的主之名所行的神迹都未能说服保禄，我们的主便急忙带着谦卑去迎接那带着骄傲的猛烈疾行往大马士革路上的保禄。这样，借着祂谦卑的话，骄傲的严酷猛烈便被阻止了。\par

40. 那位曾用谦卑的话对祂的迫害者保禄说话的主，也用谦卑的言语对法利塞人说话。因为谦卑的力量如此伟大，以致于战胜万有的天主，也未曾不带谦卑便得胜。谦卑也能在旷野中承受那顽梗之民的重担。因为对着那比众人更顽固的百姓，天主设立了那比众人更温良的梅瑟。因为一无所缺的天主，在释放了百姓之后，却需要梅瑟的谦卑，好使这谦卑能承受那激怒他的人民的愤怒和怨言。因为唯有谦卑能承受那人民的反对，连埃及的神迹和旷野中的奇事都无法制伏他们的悖逆。因为当骄傲在百姓中造成了分裂，谦卑便以祈祷弥合了他们的分裂。如果那口吃的梅瑟的谦卑能承受六十万人，那么那赐予口吃者言语的祂的谦卑，岂不更加无法估量地承受一切？因为梅瑟的谦卑，正是我们主谦卑的预像。\par

41. 主随后看到，法利塞人西满并不相信他所见的标记和奇事。主到他那里去，用谦卑的言语说服他；谦卑的话语胜过了他，而大能的奇事却没有胜过他。那么那法利塞人曾见到什么奇事呢？他曾看见死人复活，癞病人洁净，瞎子开眼。这些标记迫使那法利塞人将主当作先知款待。但那将他当作先知款待的人，转而轻蔑他，认为他无知，说（即）：\Quote{如果这人是先知，他必会知道这个妇人——那接近主的妇人——是个罪妇。} 但我们可轻视那法利塞人并说，若他是个有辨识力的人，他本应从那位接近主的罪妇学习到，祂不仅是先知，而是先知的主。因为罪妇的眼泪证明，她们所恳求的并不是一位先知，而是祂，即天主，曾因她的罪而发怒。因为先知不足以使罪人复活，先知的主降临来医治那些情况恶劣的人。但哪位医生会阻止受伤者来到他面前呢？盲目的法利塞人啊，就如那妇人来就近我们的医生那样！为什么那受伤的妇人接近祂？——她的创伤由她的泪治愈了？祂曾降临成为病人中的治疗泉源，曾宣告说：\BibleQuote{凡口渴的，让他来喝。}\BibleRef{若 7:37} 但当这人的同伴法利塞人因罪人的治愈而抱怨时，祂就自己的医术教导说，门是为病人开的，而不是为健康的人，因为\BibleQuote{健康的人不需要医生，而是有病的人。}\BibleRef{玛 9:12} 因此，医生的赞美就是病人的痊愈——为使那指责我们医生赞美的法利塞人的羞耻更大。但主惯常在街上显标记；而当祂进入法利塞人的屋子时，祂显了比在外面更大的标记。因为在街上，祂治愈了有病的身体，但在屋内，祂治愈了有病的灵魂。在外面，祂使拉匝禄复活；但在里面，祂使那罪妇复活。祂将活的灵魂还给了那已离去的尸体；并从罪妇那里驱逐了居住在她内的大罪。但那盲目的法利塞人对大事不够有能力，因为他看不见大事，就否定了所看见的小事。因为他是以色列之子，将软弱归于他的天主，而不是归于自己。因为（以色列说）：\Quote{祂虽然击打磐石，水流了出来，但祂也能给我们食粮吗？} 但当主看见他的软弱，即他错过了大事，并因大事也错过了小事，就急忙提出一个简单的言语，如同为那用奶喂养的婴孩，那婴孩还不能吃固体食物。\par

42. 因为，法利塞人啊，由你所知道的主不是先知的这一点，恰恰证明了你不认识先知。因为藉着你所说的——\Quote{如果这人是先知，他必会知道}，你在这里表明，（在你看来）凡是先知都知道一切事。但是看！有些事是对先知隐藏的；那么你怎能将一切隐藏事的启示归给先知呢？但这不明智的教师，他败坏了先知们的圣经，甚至不明白他所读的圣经。因为不仅主的伟大没有被那法利塞人辨识，他甚至没有辨识先知们的软弱。因为我们的主，作为全知者，让那个罪妇进来并接受祂的平安。但厄里叟，作为一个不知道的人，对叔能妇人说：\BibleQuote{平安归于你，\InnerQuote{和平安归于这孩子}}。\BibleRef{列下 4:26} 因此，那以为主被证明不是先知的人，他自己被证明不认识先知。当人心怀恶意而不能克制时，那恶意在其中，就诡诈地寻找开门的借口；但若那借口，即欺骗者的避难所，被驳倒时，他知道在这之内还有另一个隐藏的可以运用。\par

现在注意这个以色列人，他是如何像以色列那样顽固的。因为异教信仰被捆绑在人民的心里；因此梅瑟从他们那里被带走，好使他们内在的邪恶得以显明。但为了让他们不致蒙羞，也不让人知道他们是如何寻求偶像的，他们先寻找梅瑟，然后寻找偶像。\BibleQuote{至于这个梅瑟，我们不知道他成了什么。}\BibleRef{出 32:1} 如果那不死的天主曾领你们出埃及，为什么你们寻求一个必要死的人呢？然而他们并不想要梅瑟作他们的天主；因为梅瑟能听、能看、能责备；但他们寻求的是一个既不能听、不能看、也不能责备的天主。但一旦梅瑟死了，他身上还剩下什么呢？因为看，你们的天主是永生的天主，看，祂藉着活生生的证据将自己启示给你们。因为那时光明的云彩遮蔽他们，夜间他们有光柱。水从磐石流给他们，他们饮了它的水流。他们每天因品尝那玛纳而喜乐，我们曾听说过它的名声。梅瑟离你们有多远呢？看，梅瑟的标记环绕着你们。或者，当你们有这样的一位向导时，梅瑟这个人如何有益于你们呢？如果你们的衣服不旧，温和的空气使你们清爽，如果炎热和寒冷不伤害你们，你们从战争中得到安息，远离了来自埃及的恐惧——那么以色列还缺少什么，以致他要寻求梅瑟呢？他缺少的是公开的异教信仰。因为他寻求的不是梅瑟，而是以梅瑟不在为借口，追随了牛犊。如此我们简要地表明，当人心充满某物，但遇到反对的理由时，它就强行为之打开通向其所欲的门。\par

43. 你，法利塞人哪，渴望亵渎，你在我们的主身上看见了什么，以表明祂不是一位先知？因为看啊！属于先知们的主的事物已经在祂身上被看见了。因为涌流的眼泪急忙宣告，它们是当着天主的面流下的；悲伤的亲吻作证，它们试图赢得债权人取消债务契约；罪妇的美好香膏宣告，它是补赎的献礼。这些药物，罪妇献给了她的医师：好藉着她的眼泪，祂能洗去她的污点；藉着她的亲吻，祂能治愈她的创伤；藉着她的香膏，祂能使她的恶名如同香膏的芬芳一样甜美。这就是那位医师，祂藉着人们带给祂的药物治愈人。这些奇迹在那时被显示了；但在法利塞人眼中，代替这些出现的却是亵渎。因为在罪妇的哭泣中能建立起什么，除非祂能使罪人成义？否则，你心中审判吧，瞎眼的导师，为什么在那欢乐的筵席中有那悲伤的哭泣，以致当他们以食物欢乐时，她却在眼泪中痛苦？因为她是一个罪人，她的行为是不贞的，而这些（行为）是她惯常所做的。但如果在那时，她从罪人的放荡转向了贞洁，那么你，那个说\Quote{祂不是先知}的人，当承认祂是那使放荡者变贞洁的一位。因为你藉此知道她是一个罪人，并藉此看到她现在悔改，探究那改变她的能力在哪里。因为祂本应俯伏朝拜祂，祂默默无声，而在祂的静默中，祂使那些先知们以强烈的呼喊都不能使其转向贞洁的罪人转向了贞洁。在法利塞人的家里看见了一件奇妙且令人惊叹的事：一个罪妇坐着哭泣，而哭泣的她没有说她为什么哭泣；坐在她脚前的那位也没有对她说：\Quote{你为什么哭泣？}罪妇不需要用嘴唇向我们的主祈求，因为她相信祂，作为天主，知道隐藏在她眼泪中的祈求。我们的主也没有问她：\Quote{你做了什么？}因为祂知道她藉着纯洁的亲吻在赎她的过犯。所以，她因为相信祂知道隐藏的事，便在心中向祂献上祈祷；因为祂知道隐秘的事，不需要外在的嘴唇。那么，如果罪妇因为知道我们的主是天主，便不用嘴唇说服祂；而我们的主，因为作为天主祂辨别她的思想，所以没有质问她；你，专横的法利塞人，难道不能从两者的静默中理解两者的位置：她心中如同向天主祈祷，而祂作为天主在静默中探查她的思想吗？但法利塞人看不见、也不能理解这些事，因为他是以色列之子，虽然看见，却看不见；虽然听见，却不明白。那么，虽然我们的主知道那个法利塞人对祂怀着恶念，却温和地驳斥了他，而非严厉。因为甘甜从天降下，要打破邪恶者加在我们身上的苦毒。因此，我们的主教导那个法利塞人关于祂自己并在祂内的事，如同在说：\Quote{就像我，虽然知道你心中的恶念，却仍温和地劝导你，同样，我虽然知道这妇人的恶念，却仁慈地接纳了她。}\par

但让我们听听忍耐是怎样被急躁的思想所吸引，好把它从急躁引向理解。\BibleQuote{某个债主有两个债务人。一个欠五百德纳，另一个欠五十德纳。}——（听着，不要因比喻重复之长而厌烦，免得你与那在比喻中为施助而忍耐的主相反。）——\BibleQuote{最后，两人都无力偿还，他就免了两人的债。依你看，他们中哪一个会更爱他？\Person{西满}对祂说：\InnerQuote{我想是那多得宽免的。}主对他说：\InnerQuote{你判断得正确。}}主以祂的正义赞扬了那个悖逆的（法利塞人），因为他作出了正确的判断，虽然他在邪恶中对良善的主所行的仁慈作了回答。如今许多事蕴藏在这个比喻中；因为它是一个充满许多助益的宝库。那么，为什么主要求法利塞人为祂在两个债务人之间作判断呢？难道不是要让大的随后而来，藉着小的显示出来，好表明没有一件小事是追随大的吗？因为主既然知道隐秘的事，就忍耐着质问\Person{西满}，好使那些虽不知道却急于指责而不追究的人蒙羞。因为，人啊，在我听见你的判断之前，我并没有判断它，你为何在我没有听见有关那罪妇的案件之前，就急于指责呢？这事是为教训我们而作的，使我们快于探究而缓于宣判。因为如果那个法利塞人有忍耐，看啊！主最终赐给罪妇的宽恕，就会教导他一切。忍耐常能为获得它的人赢得一切。\par

44. 再者，藉着两个债户的宽免，我主引导那需要宽免、却视债款宽免为可憎的人进入宽免。因为法利塞人自己需要债款的宽免，但他眼中却憎恶罪妇的债款宽免。若法利塞人心中有这债款宽免，那罪妇前来向天主而非向司铎求宽免其债，就不会在他眼中视为可耻；因为司铎不能宽免那样的罪。但这罪妇因我主所行的荣耀之工，相信祂也能赦免罪过。因她知晓：谁能修复身体的肢体，也必能涤净灵魂的污点。但法利塞人虽是老师，却不明白这事。因为以色列的老师们常是愚笨的，被卑贱之人所羞辱。他们被那瞎眼的人所羞辱，他对他们说：\Quote{\Emph{我们知道这人是个罪人。}}\BibleRef{若 9:24}但他对他们说：\Quote{\Emph{祂怎样开了我的眼睛？看！天主不听罪人。}}\BibleRef{若 9:24}这些就是瞎眼的老师，成了别人的向导；而他们弯曲的路被一个瞎眼的人改正了。\par

45. 且听我主所行的奇事。因为那法利塞人以为我主不知道触摸祂的妇人是个罪人；我主便将法利塞人的嘴唇当作竖琴的弦，藉着他的嘴唇歌唱那妇人如何践踏他的罪，而他却不自知。那自以为知道而责备的人，竟成了竖琴，由另一位歌唱他所不知道的事。因为我主将罪妇的罪比作五百德纳，并藉着他所听的比喻使之传入法利塞人的耳中；又在他所下的判词中从他口中将其引出。虽然\Person{西满}下判时不知那五百德纳指的是罪妇的罪。那认为我主不知道她罪过的法利塞人，自己却当她听到比喻中的债款并以声音下判时，也发现对此毫无所知。最后当我主向他解释时，法利塞人才明白：他的耳朵和嘴唇都如同为我主所用的乐器，藉以颂扬祂知识的荣耀。\par

因为这法利塞人就是那些经师们的同伴，我主曾藉着他们自己的口定了他们的罪——\Quote{\Emph{那么，葡萄园的主人要怎样对待那些园户呢？}}\BibleRef{玛 21:40}他们对自己说：\Quote{\Emph{他要凶恶地消灭他们，并将葡萄园租给那些按时交果实的园户。}}这就是天主性，万物对祂都是容易的，祂藉着那些亵渎祂的口，甚至藉着那些口，宣布了定那些口有罪的判词。\par

46. 愿荣耀归于那不可见者，祂披上可触可见的形像，使罪人能亲近祂。因为我主没有像法利塞人所期望的那样拒绝罪妇；祂从无人能及的高天降下，正是为使卑微的税吏如\Person{匝凯}能触摸到祂。那不可触摸的本性披上了身体，正是为使万口\BibleRef{依 6:7}都能像那罪妇一样亲吻祂的脚。圣洁的灵魂藏在血肉的帐幔内，触摸所有不洁的嘴唇并圣化它们。如此，那本似因食欲而被邀请赴宴的，祂的脚却引来了眼泪；祂是良医，出来寻找那在心中寻求祂的罪妇。她，被众人践踏在尘土中的，却为我主的脚傅油，而非傅抹祂的头。因为那些自义而轻视众人的法利塞人践踏了她。但仁慈的祂，以纯净的身体圣化了她那不洁，怜视了她。\par

47. 但\Person{玛利亚}傅抹了我主身体的头，\BibleRef{玛 26:7}作为她所选的上好福分的标记。基督预言了她灵魂所选择的事。当\Person{玛尔大}为伺候忙碌时，\Person{玛利亚}渴慕从祂得到属灵事物的饱足，祂也以身体事物饱足我们。所以\Person{玛利亚}以珍贵的香液使祂舒畅，如同祂以崇高的教导使她舒畅。\Person{玛利亚}以油彰显了祂可朽的奥迹，祂以教导杀死了她肉体的私欲。如此，罪妇以泪水的洪流，满有信心地在祂脚下获得了罪赦；那患血漏的妇人，从祂衣边偷得了医治。但\Person{玛利亚}公开从祂口中得到了祝福，作为她亲手事奉祂头的报酬。因为她将珍贵的香液倒在祂头上，从祂口中得到了奇妙的应许。这就是那播在上面、结在下面的香液。她将其播在祂头上，从祂唇间收取果实——\Quote{\Emph{无论在什么地方，这福音被传扬，她也要被传扬，作为纪念。}}\BibleRef{玛 26:13}因此，她当时从祂所领受的，祂能使之传至万代；任何世代都不能阻止。因为那她倒在祂头上的香液，在众宾客前散发芬芳，使祂舒畅；同样，祂赐予她的美名传至万代，给她带来荣耀。正如所有赴宴的人都感知到她的香液；所有来到世界的人也应该感知到她的胜利。这是在各世代都被要求增息的贷款。\par

48. 那司铎 \Person{西默盎} 在怀中抱起了祂，为将祂呈献于天主面前 \BibleRef{路 2:28}，他看着祂，明白了他不是在呈献祂，而是自己被呈献。因为圣子并非由仆人呈献给圣父，而是仆人经由圣子呈献给祂的主。因为那藉着祂一切祭献得以呈献的，不可能由他人呈献。因为祭献并不呈献那献祭者，而是献祭者呈献祭品。因此，那接受祭献的，将自己交予他人呈献，好使那些呈献祂的人，在献祭之时，也经由祂被呈献。正如祂将自己的身体赐下作为食物，使吃祂的人因吃而得到生命，同样，祂将自己交予献祭，好使那些献祭祂的人之手，因祂的十字架而被圣化。因此，尽管 \Person{西默盎} 的双臂看似在呈献圣子，但 \Person{西默盎} 的话却证实他是被圣子所呈献。因此，我们对此无可争辩，因为所说的话终止了争辩——\BibleQuote{主啊！现在请祢的仆人平安离去吧。}\BibleRef{路 2:29} 那被允许平安离去归于天主的人，便是作为祭品呈献于天主。而为了显明他是经由谁被呈献，他说：\BibleQuote{因为我亲眼看见了祢的仁慈。}\BibleRef{路 2:30} 若他身上没有恩宠成就，为何他感恩呢？然而他感恩正当合理，因为他被认为配得用双臂怀抱那位天使和先知们极度渴望看见的。因为：\InnerQuote{我亲眼看见了祢的仁慈。} 那么，让我们明白并看清。\Quote{仁慈}是那向他人施予仁慈的，还是那从他人领受仁慈的？但若仁慈是那向众人施予仁慈的，那么 \Person{西默盎} 便合宜地以施予他仁慈的仁慈之名来称呼我们的主——祂使他脱离那满是陷阱的世界，得以前往那充满喜乐的 \Place{伊甸}。因为那司铎说了并证实自己作为祭品被呈献，好叫他远离这败坏的世界，被存留在那稳妥的宝库中。因为一个人若有可能失去他所获得的，他就应谨慎使之得以保全。但我们的主不可能失去什么，反而是经由祂，失去的被寻回。因此，藉着那不能失去的圣子，那极其不愿失去的仆人被呈献了。\BibleQuote{我亲眼看见了祢的仁慈。} 显然 \Person{西默盎} 从他怀抱的那婴孩领受了恩宠。因为他从内里从那婴孩领受了恩宠，虽然他外在用双臂怀抱祂。因为经由那荣耀的，甚至当他被抱着时弱小卑微，那抱着他的却变得伟大。\par

49. 但既然 \Person{西默盎} 忍受以他软弱的双臂怀抱那受造物无法承受的威严，显然他的软弱是被他所怀抱的力量所加强。因为那时 \Person{西默盎} 与所有受造物一道，被圣子的全能力量暗中支撑。这真是奇妙：外表上，是那被加强的抱着那加强他的；但内里，是那力量托着那托着它的。因为那威严自我约束，好使抱祂的人能承受祂；为使那威严屈尊俯就我们的卑微，我们的爱情也应从一切欲望中被提升，以达到那威严。\par

50. 同样，那艘载着我们主的船：是祂托着它，因祂制止了风，以免它沉没。\Quote{安静，因为你们被困住了。}当祂在海上时，祂的手臂甚至达到了风的源头 \BibleRef{谷 4:39}，将它封住。那船托着祂的人性，但祂天主性的能力托着船和船上的一切。但为显示祂的人性本不需要船，祂像创造万物的建筑师一样，不用船匠拼合固定的木板，而是使水坚固、结合，置于祂脚下。同样，主坚固了司铎 \Person{西默盎} 的双手，使他的双臂能在圣殿中托起那托起万有的力量；正如祂坚固了宗徒 \Person{西默盎} 的双足，使它们能在水上托起自己。因此，那名号先在圣殿中托着首生者 \OriginalTerm{首生者}{first-begotten}，后来在海中被首生者托起；为显示如同在海中那溺水者被祂托起一样，在旱地上祂并不需要被 \Person{西默盎} 托着。但我们的主公开地在海中托起了 \Person{西默盎}，以教导我们，即使在旱地上，祂也暗中扶持着他。\par

51. 因此，圣子来到仆人面前，并非要由仆人将圣子呈献，而是借圣子使仆人能向祂的主呈献司祭职和先知职，与祂一同保存。因为藉梅瑟所赐下的先知职和司祭职，这两者都得以传承，一直传到西默盎。他原是一个纯洁的器皿，他圣化了自己，为能像梅瑟一样，堪当两者。有些小器皿能承载大恩赐。有些恩赐，由于它们的恩宠，一人能够堪当；然而许多人因它们的伟大而不能堪当。于是，西默盎呈献了我们的主，并在祂身上献上了这两样；如此，在旷野中赐给梅瑟的，在圣殿中从西默盎那里领受了。但由于我们的主是那器皿，里面充满一切丰盛，当西默盎在圣殿中向天主献上祂时，他（作为奠祭）将这两样（恩赐）倾注在祂身上：司祭职出自祂的双手，先知职出自祂的口唇。司祭职因西默盎的洁净而持续在他的双手上；先知职因启示而在他的口唇上运行。当这两种能力看见了那两者之主时，它们二者联合起来，倾注到那能承载两者的器皿中；它能容纳司祭职、王权和先知职。那婴孩，因祂的慈爱而被裹在襁褓中，却因祂的尊威而穿上司祭职和先知职。因为西默盎将这两样穿在祂身上，并将祂交给那位曾用襁褓包裹祂的人。当他将祂交给祂的母亲时，他同时交出了司祭职；当他向母亲论及祂而预言说：\BibleQuote{这孩子被立为堕落和复兴} \BibleRef{路 2:34} ，他也将先知职一同交给了祂。\par

52. 于是玛利亚接过她的长子，便出去了。祂外表被裹在襁褓中，但暗地里却穿上了先知职和司祭职。凡从梅瑟传承下来的，都从西默盎那里领受了，但继续存在并被那两者之主所拥有。于是，先是管家，后是库房管理员，将司祭职和先知职的钥匙交给了那对两者的库房有权柄者。因此，祂的父将圣神赐给祂，并不限量，\BibleRef{若 3:34} 因为圣神的一切量度都在祂手中。为显示祂从先前的管家那里领受了钥匙，我们的主对西默盎说：\BibleQuote{我要将门的钥匙交给你} \BibleRef{玛 16:19} 。但祂若未曾从他人那里领受，又怎能将钥匙交给他人呢？所以，祂从司祭西默盎那里领受的钥匙，又交给了另一位宗徒西默盎；这样，即使百姓未曾听从前一位西默盎，外邦人却要听从后一位西默盎。\par

53. 但因为若翰也是圣洗的库房管理员，那管家的主人来到他那里，从他那里领受和好之家的钥匙。因为若翰惯常以普通的水洗去罪过的污点，使身体堪当穿上由我们的主所赐的圣神之衣。因此，因为圣神与圣子同在，祂来到若翰那里，从他那里领受圣洗，为将无形的圣神与有形的水混合；这样，那些身体感受到水的浸润的人，他们的灵魂也能感受到圣神的恩赐；就如身体外表感受到水倾注在他们身上，同样灵魂内在感受到圣神的倾注。因此，正如我们的主受洗时，穿上了圣洗并带着圣洗，同样，当祂在圣殿中被呈献时，祂穿上了先知职和司祭职，带着司祭职的纯洁在祂纯洁的肢体上，带着先知的话语在祂奇妙的耳朵里。因为当西默盎圣化那圣化万物的孩子的身体时，那身体在其圣化中接受了司祭职。再者，当西默盎论及祂而说预言时，先知职迅速进入孩子的听觉。因为若翰曾在母胎中跳跃，并感知到我们主之母的声音，\BibleRef{路 1:41} 那么我们的主在圣殿中岂不是更应该听到吗？看！正是因为祂，若翰才得以在母胎中知道<听>。\par

54. 因此，每一个为圣子所储存的恩赐，祂都从它们真正的树上收集而来。因为祂从约旦河领受了圣洗，尽管若翰在祂之后仍施行圣洗。祂从圣殿领受了司祭职，尽管大司祭亚纳斯仍在执行。再者，祂领受了在义人中所传承的先知职，尽管盖法用它轻蔑地为我们的主编了一个冠冕；祂从达味家领受了王权，尽管黑落德占据其位并执行之。\par

55. 这就是那飞翔并自高天降下的一位；当祂昔日赐给古人的所有恩赐看见祂时，它们从各方飞来，停留在它们的赐予者身上。因为它们从四面聚集，来嫁接回它们的本树。因为它们曾被嫁接到苦树上，即邪恶的君王和司祭。因此它们急忙来到它们甘甜的母本，即天主性，祂丰丰富富地降下到以色列人民中间，为叫祂的各部分都能聚集到祂身上。当祂从它们中领受了祂自己的，那不属于祂的就被拒绝了；因为为了祂自己的缘故，祂也容忍了那不属祂的。因为祂容忍了以色列的偶像崇拜，为了祂的司祭职；祂容忍了他们的占卜者，为了祂的先知；祂容忍了他们的邪恶政权，为了祂神圣的冠冕。\par

56. 但当我们的主从他们那里将司祭职归于自己时，祂借此圣化了所有外邦人。再者，当祂将先知职归于自己时，祂借此向万民启示了祂的计谋。当祂编织祂的冠冕时，祂捆绑了那俘掳所有人的强者，并分取了他的掠物。这些恩赐与无花果树一样荒芜，那无花果树虽不结果，却使如此光荣的能力变得荒芜。因此，因它不结果子，它被砍下，好让这些恩赐从它而出，并在所有外邦人中结出丰硕的果实。\par

57. 因此，那位来使我们身体成为祂内住之住所的祂，越过了所有那些居所。那么我们每个人都要成为爱我的那一位的居所。让我们到祂那里去，与祂同住。这就是天主性，虽然整个受造界不能容纳祂，但一个谦卑朴实的灵魂却足以接纳祂。\par

\SourceRecord{Translated by A. Edward Johnston.；Nicene and Post-Nicene Fathers, Second Series；Vol. 13.；Edited by Philip Schaff and Henry Wace.；Buffalo, NY: Christian Literature Publishing Co.,；1898.；<http://www.newadvent.org/fathers/3706.htm>.}
